\section{Related Work}
\label{sec:related-work}

\subsection{Attention Optimizations}

Multi-Head Attention (MHA)~\citep{vaswani2017transformer} faces computational and IO challenges. FasterTransformer~\citep{fastertransformer} reduces global memory footprint via Fused Multi-Head Attention (FMHA), but doesn't scale to long contexts because shared memory usage is linear to sequence length.
ByteTransformer~\citep{zhai2023bytetransformer} optimizes FMHA on variable-length input.
FlashAttention ~\citep{dao2022flashattention} uses online-softmax ~\citep{milakov2018onlinesoftmax} trick to reduce the shared memory footprint to constant size, enabling long contexts. FlashAttention2\&3~\citep{dao2023flashattention2,jay2024flashattention3} further optimizes FlashAttention by improving loop structure and overlapping softmax and GEMM. FlashDecoding ~\cite{dao2023flash-decoding} applies Split-K to decode attention kernels. LeanAttention~\cite{sanovar2024lean} uses StreamK~\citep{osama2023streamk} to reduce wave-quantization~\citep{wave-quantization} in attetnion (with fixed sequence length). \sys extends the FlashAttention2\&3 template to support sparse attention kernels, while using StreamK-like optimizations on variable length sequences. Nanoflow ~\citep{zhu2024nanoflow} introduces horizontal fusion of GEMM, attention, and communication operations, while POD-Attention ~\citep{kamath2024podattention} focuses on optimizing chunked-prefill attention. The JIT compilation framework of \sys can be extended to generate kernels supporting these fusion techniques. FlashDecoding++~\citep{flashdecodingpp} leverages attention scale statistics to predefine a unified max value. This process converts attention composition (section \ref{sec:attention-composition}) to summation, enabling TMA Store Reduce ~\citep{tma-store-reduce} to asynchronously updating global \textit{attention state}s, it's orthogonal to \sys's contribution and we leave it for future work.

Recent works like RelayAttention~\citep{zhu2024relayattention}, Hydragen~\citep{juravsky2024hydra}, ChunkAttention~\citep{ye2024chunkattention}, and Parrot~\citep{lin2024parrot} explore shared prefix decoding attention but require separate KV-Cache management for prefixes and suffixes. In contrast, \sys's composable formats support multi-level, multiple-prefix decoding with unified page table management, enabling seamless integration into LLM serving frameworks without modifying memory management modules.

\subsection{Sparse Optimizations on GPUs}

FusedMM~\citep{rahman2021fusedmm} explores Sparse-dense Matrix Multiplication (SpMM) fusion, though it omits softmax computation, limiting direct applicability for accelerating attention. \citet{zhang2022fusedgat} explore Graph Attention Networks (GAT) kernel fusion, SAR~\citep{mostafa2022sar} serializes Sparse Attention aggregation, akin to FlashAttention, neither work explores using Tensor Cores. Blocksparse library~\citep{gray2017gpu-bsr} implements BSR GEMM with tensor cores. \citet{chen2021tensorcoressparsity}, TC-GNN~\citep{wang2023tcgnn} and Magicube~\citep{li2022magicube} propose vector sparse formats to leverage Tensor Cores effectively. \sys improves upon these to support any block sizes $(b_r, b_c)$ in FlashAttention.

\subsection{Attention Compilers}

FlexAttention~\citep{he2024flexattention} provides a user-friendly interface for programming attention variants, compiling them into block-sparse flashattention implemented in Triton~\citep{triton}. It uses PyTorch Compiler~\citep{jason2024pytorch2} to automatically generate backward passes. \sys expands the FlexAttention's programming interface to support query/key transformations, and focus on vector-sparsity and load-balancing for LLM serving. \sys generates CUDA code instead of Triton because Triton still underperform CUDA \& CUTLASS in many use cases. \sys can act as a backend for FlexAttention in forward pass. Mirage~\citep{wu2024mirage} optimizes tiling strategies for GEMM and FlashAttention using a probabilistic equivalence verifier, relying on Triton and CUTLASS for code generation. However, it lacks support for variable length and sparse data structures, and doesn't include safe-softmax, unlike \sys, which is directly applicable to LLM serving.

\subsection{LLM Serving Systems}

Orca~\citep{yu2022orca} introduces continuous batching for enhanced throughput. PagedAttention~\citep{kwon2023vllm} uses a Page Table for KV-Cache management. Sarathi-serve~\citep{agrawal2024sarathi} improves efficiency by piggybacking decode operations with chunked-prefill, while SGLang~\citep{zheng2023sglang} utilizes RadixTree for better prefix-caching and KV-management. \sys provides a unified solution for these attention mechanisms through block-sparse attention kernels. vAttention~\citep{prabhu2024vattention} shows that GPU virtual memory can manage address translation in PageAttention without special kernels. Yet, challenges like dynamic KV-Cache sparsity persist, as seen in Quest~\citep{tang2024quest}. Here, \sys's block sparse kernel remains effective. Additionally, \sys can be combined with vAttention by generating kernels for contiguous KV-Cache storage.
